\documentclass{article}

\usepackage[spanish,english]{babel}
\addto\shorthandsspanish{\spanishdeactivate{~<>}}
\usepackage[utf8]{inputenc}
\usepackage[T1]{fontenc}
\usepackage{lmodern}
\usepackage{microtype}
\usepackage{graphicx}
\usepackage[margin=25mm,inner=30mm]{geometry}
\usepackage{fancyhdr}
\usepackage{listings}
\usepackage[colorlinks=true,linkcolor=blue,urlcolor=blue]{hyperref}
\graphicspath{{recursos/}}

\newcommand{\theCampus}{Campus Vitacura}
\newcommand{\theAssignment}{Sistemas de Información para la Gestión}
\newcommand{\theReportTitle}{Laboratorio 5: SQL}
\newcommand{\theGroup}{Entrega individual · ICN-292 · Paralelo 100}
\newcommand{\theAuthorA}{Joaquín Tapia Mandiola}
\newcommand{\theCareer}{Ingeniería Civil Industrial}
\newcommand{\theRepository}{https://github.com/JotaReyyyyyyyyyyyy/ICN292-Lab5-Tapia-Joaquin}

\setlength{\headheight}{45pt}
\setlength{\headsep}{12pt}
\pagestyle{fancy}
\fancyhf{}
\lhead{\includegraphics[scale=0.3]{Logo usmUSM.png}}
\renewcommand{\headrulewidth}{0.5pt}
\fancyfoot[C]{\textbf{\thepage}}

\lstset{
  language=SQL,
  basicstyle=\ttfamily\footnotesize,
  columns=fullflexible,
  keepspaces=true,
  upquote=true,
  showstringspaces=false,
  breaklines=false,
  aboveskip=3pt,
  belowskip=4pt,
  xleftmargin=6pt,
  xrightmargin=6pt,
  literate={ñ}{{\~n}}1 {á}{{\'a}}1 {é}{{\'e}}1 {í}{{\'i}}1 {ó}{{\'o}}1 {ú}{{\'u}}1
}
\newcommand{\evidencia}[2]{%
  \par\noindent\includegraphics[width=#2]{consulta-#1.png}\par
}
\hypersetup{pdftitle={ICN-292 - Laboratorio 5: SQL},pdfauthor={Joaquín Tapia Mandiola},pdfcreator={}}

\begin{document}
\selectlanguage{spanish}
\hypersetup{pageanchor=false}
\begin{titlepage}
\thispagestyle{empty}
\vspace*{1.2cm}
\begin{center}
\includegraphics[scale=0.5]{Logo usmUSM.png}
\end{center}
\vspace{2.2cm}
\begin{center}
\hrule
\vspace{2mm}
{\Huge\bfseries \theReportTitle}\\[0.5cm]
{\Large\scshape \theAssignment}
\vspace{5mm}
\hrule
\vspace{1cm}

\textbf{\theGroup}\\
\vspace{0.8cm}
\textbf{\theAuthorA}\\
{\small \theCareer}\\[0.7cm]
{\small Repositorio GitHub}\\[0.2cm]
{\small \expandafter\url\expandafter{\theRepository}}

\vfill
\textsc{Noviembre, 2026}\\[0.5cm]
\textsc{\theCampus}
\end{center}
\end{titlepage}
\hypersetup{pageanchor=true}
\pagenumbering{arabic}

\section*{Consulta 1}
\begin{lstlisting}
SELECT v.nombre, COUNT(ve.id_venta) AS cantidad_ventas
FROM vendedores AS v
LEFT JOIN ventas AS ve ON ve.id_vendedor = v.id_vendedor
GROUP BY v.id_vendedor, v.nombre
ORDER BY cantidad_ventas DESC, v.nombre;
\end{lstlisting}
\evidencia{01}{74mm}
\vspace{5mm}
\section*{Consulta 2}
\begin{lstlisting}
SELECT t.nombre AS tienda,
       COUNT(ve.id_venta) AS cantidad_ventas,
       COALESCE(SUM(ve.unidades), 0) AS unidades_vendidas
FROM tiendas AS t
LEFT JOIN vendedores AS v ON v.id_tienda = t.id_tienda
LEFT JOIN ventas AS ve ON ve.id_vendedor = v.id_vendedor
GROUP BY t.id_tienda, t.nombre
ORDER BY unidades_vendidas DESC, t.nombre;
\end{lstlisting}
\evidencia{02}{88mm}
\clearpage
\section*{Consulta 3}
\begin{lstlisting}
SELECT t.nombre, t.comuna
FROM tiendas AS t
LEFT JOIN vendedores AS v ON v.id_tienda = t.id_tienda
WHERE v.id_vendedor IS NULL;
\end{lstlisting}
\evidencia{03}{74mm}
\vspace{5mm}
\section*{Consulta 4}
\begin{lstlisting}
SELECT b.modelo
FROM bicicletas AS b
LEFT JOIN ventas AS ve
       ON ve.id_bicicleta = b.id_bicicleta
      AND ve.fecha >= '2025-06-01' AND ve.fecha < '2025-07-01'
WHERE ve.id_venta IS NULL
ORDER BY b.modelo;
\end{lstlisting}
\noindent\footnotesize La condición del mes va en el \texttt{ON} del \texttt{LEFT JOIN}; si la hubiera puesto en el \texttt{WHERE}, la consulta no devolvería filas, porque los modelos sin venta en junio llegan con la fecha en \texttt{NULL} y no cumplen el filtro.\par\normalsize
\smallskip
\evidencia{04}{74mm}
\clearpage
\section*{Consulta 5}
\begin{lstlisting}
SELECT v.nombre, t.nombre AS tienda
FROM vendedores AS v
LEFT JOIN tiendas AS t ON t.id_tienda = v.id_tienda
LEFT JOIN ventas AS ve ON ve.id_vendedor = v.id_vendedor
WHERE ve.id_venta IS NULL
ORDER BY v.nombre;
\end{lstlisting}
\evidencia{05}{74mm}
\vspace{5mm}
\section*{Consulta 6}
\begin{lstlisting}
SELECT a.nombre AS vendedor_antiguo, a.fecha_ingreso AS ingreso_antiguo,
       b.nombre AS vendedor_nuevo, b.fecha_ingreso AS ingreso_nuevo,
       t.nombre AS tienda
FROM vendedores AS a
JOIN vendedores AS b
  ON b.id_tienda = a.id_tienda
 AND (a.fecha_ingreso < b.fecha_ingreso
      OR (a.fecha_ingreso = b.fecha_ingreso AND a.id_vendedor < b.id_vendedor))
JOIN tiendas AS t ON t.id_tienda = a.id_tienda
ORDER BY t.id_tienda, a.fecha_ingreso, b.fecha_ingreso;
\end{lstlisting}
\evidencia{06}{135mm}
\clearpage
\section*{Consulta 7}
\begin{lstlisting}
SELECT t.nombre AS tienda, SUM(ve.unidades * b.precio) AS monto_montana_ruta
FROM tiendas AS t
JOIN vendedores AS v ON v.id_tienda = t.id_tienda
JOIN ventas AS ve ON ve.id_vendedor = v.id_vendedor
JOIN bicicletas AS b ON b.id_bicicleta = ve.id_bicicleta
WHERE b.categoria IN ('Montaña', 'Ruta')
GROUP BY t.id_tienda, t.nombre
HAVING SUM(ve.unidades * b.precio) > 2000000
ORDER BY monto_montana_ruta DESC;
\end{lstlisting}
\evidencia{07}{74mm}
\vspace{5mm}
\section*{Consulta 8}
\begin{lstlisting}
SELECT modelo, precio
FROM bicicletas
WHERE precio > (SELECT AVG(precio) FROM bicicletas)
ORDER BY precio DESC;
\end{lstlisting}
\evidencia{08a}{74mm}
\smallskip
\begin{lstlisting}
SELECT AVG(precio) AS promedio_catalogo
FROM bicicletas;
\end{lstlisting}
\evidencia{08b}{74mm}
\smallskip
\noindent\footnotesize El promedio del catálogo es 657.490,00. Lo calcula la subconsulta \texttt{SELECT AVG(precio) FROM bicicletas} del \texttt{WHERE}; la segunda sentencia lo muestra por separado.\par\normalsize
\clearpage
\section*{Consulta 9}
\begin{lstlisting}
SELECT v.nombre, AVG(ve.unidades) AS promedio_unidades
FROM vendedores AS v
JOIN ventas AS ve ON ve.id_vendedor = v.id_vendedor
GROUP BY v.id_vendedor, v.nombre
HAVING AVG(ve.unidades) > (SELECT AVG(unidades) FROM ventas)
ORDER BY promedio_unidades DESC, v.nombre;
\end{lstlisting}
\evidencia{09}{74mm}
\vspace{5mm}
\section*{Consulta 10}
\begin{lstlisting}
SELECT v.nombre, SUM(ve.unidades * b.precio) AS monto_total
FROM vendedores AS v
LEFT JOIN ventas AS ve ON ve.id_vendedor = v.id_vendedor
LEFT JOIN bicicletas AS b ON b.id_bicicleta = ve.id_bicicleta
GROUP BY v.id_vendedor, v.nombre
ORDER BY monto_total DESC;
\end{lstlisting}
\noindent\footnotesize Fernanda Lillo y Rubén Castro muestran \texttt{NULL} porque no tienen ventas: el \texttt{LEFT JOIN} los deja sin filas que sumar y \texttt{SUM} devuelve \texttt{NULL}. Para que mostraran cero habría que usar \texttt{COALESCE(SUM(ve.unidades * b.precio), 0)}.\par\normalsize
\smallskip
\evidencia{10}{74mm}
\end{document}
